\chapter{Closed three-manifold topology and the Poincare problem}
\label{ch:foundations-v4}

Simple connectedness gives a closed smooth three-manifold the homotopy
type of a sphere. The Ricci-flow argument will strengthen this to a
diffeomorphism. This chapter constructs the initial topology needed by
surgery and extinction, without assuming a sphere recognition theorem.

\section{Manifolds and the two conclusions}
\label{sec:foundations-endpoints}

\begin{definition}[The smooth and topological problems]
\label{def:foundations-manifold}
The smooth Poincare proposition quantifies over compact, Hausdorff,
second-countable, simply connected spaces with a smooth atlas over
$\mathbb R^3$, and asks for a diffeomorphism with the standard unit
sphere $S^3\subset\mathbb R^4$. The topological proposition instead
requires only a topological atlas and asks for a homeomorphism.
Here simple connectedness includes nonemptiness and path connectedness;
the manifold charts have no boundary, and smooth means $C^\infty$.
\lean{PoincareMT.SmoothPoincare}
\lean{PoincareMT.TopologicalPoincare}
\end{definition}

\begin{theorem}[Smooth and topological Poincare conclusions]
\label{thm:foundations-targets}
Both propositions of Definition~\ref{def:foundations-manifold} hold.
\lean{PoincareMT.smoothPoincareSkeleton}
\lean{PoincareMT.topologicalPoincareSkeleton}
\uses{def:foundations-manifold}
\source{MT2007}
\dcref{Corollary 0.2(a), p. xii, for the smooth conclusion}
\end{theorem}

These are the book's endpoints, proved in Chapters 14 and 18. In this
chapter $M$ denotes a closed, simply connected smooth three-manifold.
Homology and cohomology have integer coefficients, $\simeq$ denotes
homotopy equivalence, and $\cong$ denotes homeomorphism.

\begin{proposition}[Initial topology]
\label{prop:foundations-initial}
The chosen atlas of $M$ has compatible orientation signs. There is a
finite CW structure of dimension at most three on $M$, and for every
$p\in M$,
\[
 \pi_1(M,p)=0,\qquad \pi_2(M,p)=0,\qquad
 \pi_3(M,p)\cong\mathbb Z.
\]
Moreover, $M\simeq S^3$.
\lean{PoincareMT.closedSimplyConnectedThreeManifoldTopology}
\uses{prop:foundations-triangulation,lem:foundations-determinant-cover,cor:foundations-homotopy-groups,prop:foundations-cw-equivalence}
\source{MT2007}
\dcref{discussion after Proposition 18.9, p. 424}
\dcref{discussion after Claim 18.20, p. 431}
\end{proposition}

The following constructions prove this proposition. The nonzero third
homotopy class supplies the width argument; orientation supplies the
obstruction to a two-sided projective plane.

\section{A finite three-dimensional CW structure}
\label{sec:foundations-triangulation}

The construction embeds $M$ in Euclidean space, perturbs an ambient grid
away from low-dimensional faces, and projects an associated order complex
back to $M$. Its quantitative steps are recorded separately.

\begin{lemma}[Finite affine-plane avoidance]
\label{lem:foundations-affine-avoidance}
For $N>0$, a ball $B(c,R)\subset\mathbb R^N$ with $R>0$, and
$m$ nonempty proper affine subspaces $P_i$, there is $x\in B(c,R)$
such that
\[
 \operatorname{dist}(x,P_i)>
 \frac{R}{4(N+1)^2(m+1)}\qquad(1\leq i\leq m).
\]
\lean{Poincare.Topology.exists_ball_point_avoiding_affine_planes}
\end{lemma}
\begin{proof}
Use the $(m+1)^N$ points with coordinates
$c_j+\frac{R}{N+1}(\frac{2a_j+1}{m+1}-1)$,
$0\leq a_j\leq m$. A unit normal to each forbidden hyperplane has
one coordinate of absolute value at least $1/(N+1)$. Thus at most
$(m+1)^{N-1}$ grid points lie in its prescribed slab. The $m$ slabs
cannot cover the grid.
\end{proof}

\begin{lemma}[Clearance across a join]
\label{lem:foundations-join-clearance}
Let $P\subset P^*$ be nonempty affine subspaces and $Q$ a nonempty
finite subset of $P^*$. Suppose $d,t,D>0$, the convex hull of $Q$
has distance at least $d$ from $P$, $z$ has distance at least $t$ from
$P^*$, and $\|z-y\|\leq D$ on that hull. Then
\[
 \operatorname{dist}(x,P)\geq\frac{dt}{D+t}
 \quad(x\in\operatorname{conv}(Q\cup\{z\})).
\]
\lean{Poincare.Topology.affine_join_distance_lower_bound}
\end{lemma}
\begin{proof}
At $x=(1-a)y+az$, orthogonal projection gives lower bounds $at$ and
$d-aD$. Their maximum is at least $dt/(D+t)$.
\end{proof}

\begin{proposition}[Triangulation of an embedded three-manifold]
\label{prop:foundations-embedded-triangulation}
Let $M$ be a nonempty compact Hausdorff smooth three-manifold and let
$e:M\to\mathbb R^N$, $N>3$, be a smooth closed embedding with
injective derivative. There is a finite poset $I$ and a homeomorphism
$|\operatorname{Ord}(I)|\cong M$; every chain in $I$ has at most
four elements.
\lean{Poincare.Topology.exists_embedded_three_finite_triangulation}
\uses{lem:foundations-affine-avoidance,lem:foundations-join-clearance}
\end{proposition}
\begin{proof}
Uniform graph charts and a tubular neighborhood follow from compactness.
Perturb a mesh-$h$ grid using the clearance lemmas. Faces of dimension
at most $N-4$ stay a distance $bh>0$ from $e(M)$, with simplex
diameters and altitudes controlled by fixed multiples of $h$.

Call a face active when it meets $e(M)$. Normal projection in graph
charts shows that minimal active faces have $N-2$ vertices and meet
$e(M)$ at unique fixed points of contractions. Average these points to
obtain centers with positive barycentric coordinates in every active face.

Order active faces by inclusion. Their cardinalities range from $N-2$
to $N+1$, so flags have at most four terms. Mapping flag vertices to
centers gives an injective affine realization: successive largest supports
recover each barycentric coefficient uniquely.

Compare it with the tangent-plane flag realization. The clearance and
altitude bounds make the displacement Lipschitz constant at most $1/2$
when the graph error is sufficiently small. The normal-fiber equation
is then a contraction, giving exactly one point over each point of
$e(M)$. Nearest-point projection is a continuous bijection from the
compact flag realization to $e(M)$, hence a homeomorphism.
\end{proof}

\begin{proposition}[Finite triangulated model]
\label{prop:foundations-triangulation}
Every nonempty compact Hausdorff smooth three-manifold is homeomorphic
to the order complex of a finite poset whose chains have at most four
elements. It has a finite CW structure of dimension at most three.
\lean{Poincare.Topology.exists_compact_three_manifold_finite_triangulation}
\lean{Poincare.Topology.exists_finite_simplicialCW}
\lean{Poincare.Topology.exists_finiteCW_of_homeomorph}
\uses{prop:foundations-embedded-triangulation}
\end{proposition}
\begin{proof}
Take Whitney's compact smooth embedding, append four zero coordinates,
and apply the embedded construction. Simplex interiors give a finite
CW structure on the realization. Transport its characteristic maps and
weak topology along the homeomorphism to $M$.
\end{proof}

\section{Two orientation constructions}
\label{sec:foundations-orientation}

\begin{definition}[Compatible chart signs]
\label{def:foundations-chart-signs}
For each chart $x_i:U_i\to\mathbb R^3$, choose a locally constant
$\varepsilon_i:U_i\to\{1,-1\}$. Compatibility means
\[
 \varepsilon_i(p)\varepsilon_j(p)
 \det D(x_j\circ x_i^{-1})_{x_i(p)}>0
 \qquad(p\in U_i\cap U_j).
\]
\lean{PoincareMT.OrientationCompatibleAtlas}
\uses{def:foundations-manifold}
\end{definition}

\begin{lemma}[The determinant orientation cover]
\label{lem:foundations-determinant-cover}
A simply connected smooth three-manifold admits compatible chart signs.
\lean{Poincare.Topology.nonempty_orientationCompatibleAtlas}
\uses{def:foundations-chart-signs}
\source{Hatcher}\dcref{Proposition 1.33, p. 61}
\end{lemma}
\begin{proof}
Transition determinant signs define a two-sheeted cover. Lift the
identity after choosing a point over a basepoint. Simple connectedness
makes the lift a section, whose coordinates give the compatible signs.
\end{proof}

\label{def:foundations-local-orientation}
For $p\in M$, write $L_p=H_3(M,M\setminus\{p\})$.
A local integral orientation chooses generators $\omega_p\in L_p$
locally represented by a common class in $H_3(M,M\setminus U)$ on
a neighborhood $U$. The next theorem explicitly constructs this data.

\begin{lemma}[Coherent generators]
\label{lem:foundations-coherent}
A simply connected topological three-manifold admits a local integral
orientation. No smooth structure is required.
\lean{Poincare.Topology.exists_integralThreeLocallyRepresentedGenerators}
\uses{def:foundations-manifold}
\source{Hatcher}\dcref{Section 3.3, pp. 233--235}
\end{lemma}
\begin{proof}
Excision identifies local homology groups with $\mathbb Z$. Their two
generators form a covering, since changes of local frame act by locally
constant signs. Lift the identity to obtain the coherent generators and
their local representatives.
\end{proof}

\section{The third homology group}
\label{sec:foundations-top-homology}

\begin{lemma}[Compact-support detection and gluing]
\label{lem:foundations-support-gluing}
On a Hausdorff topological three-manifold, a class in
$H_3(M,M\setminus K)$, $K$ compact, is determined by its restrictions
to local groups at points of $K$. A locally represented orientation
has a unique such class with the prescribed restrictions.
\lean{Poincare.Topology.integralThreeManifoldCompactSupport}
\lean{Poincare.Topology.exists_unique_integralSupportHomology_compact_gluing}
\uses{lem:foundations-coherent}
\source{Hatcher}\dcref{Lemma 3.27, pp. 236--238}
\end{lemma}
\begin{proof}
The Euclidean relative calculation gives detection and vanishing above
degree three on small compact supports. Relative Mayer--Vietoris extends
them to finite unions. Finite chain supports and a finite chart cover
reduce general compact supports to these cases. Exactness glues the local
classes, and detection gives uniqueness.
\end{proof}

\begin{proposition}[Top integral homology]
\label{prop:foundations-top-homology}
If $M$ is a closed connected topological three-manifold with local
integral orientation $\omega$, there is a unique class $[M]$ restricting
to $\omega_p$ at every point, and $z\mapsto z[M]$ is an isomorphism
$\mathbb Z\to H_3(M)$.
\lean{Poincare.Topology.nonempty_integralSupportHomologyUniv_equiv_int}
\uses{lem:foundations-support-gluing}
\source{Hatcher}\dcref{Theorem 3.26, p. 236}
\end{proposition}
\begin{proof}
Glue with $K=M$. The coefficient of any class relative to $\omega_p$
is locally constant, hence constant. Detection makes the coefficient
map injective, and $[M]$ maps to $1$.
\end{proof}

\section{Cap products and vanishing second homology}
\label{sec:foundations-second-homology}

\begin{definition}[The open-set cap property]
\label{def:foundations-cap-property}
Fix compact-support detection and a locally represented orientation.
For open $U\subset M$, let $D_U^q:H_c^q(U)\to H_{3-q}(U)$ be
cap product with the orientation. The cap property consists of
\begin{equation}
\label{eq:foundations-cap-property}
\begin{aligned}
 &D_U^1\text{ is surjective},\\
 &D_U^2\text{ and }D_U^3\text{ are isomorphisms},\\
 &H_c^q(U)=0\quad(q\geq4).
\end{aligned}
\end{equation}
\lean{Poincare.Topology.integralOpenCapProperty}
\uses{lem:foundations-support-gluing}
\end{definition}

\begin{lemma}[Union step]
\label{lem:foundations-cap-union}
If $U,V$, and $U\cap V$ satisfy the cap property, then so does $U\cup V$.
\lean{Poincare.Topology.integralOpenCapProperty_union}
\uses{def:foundations-cap-property}
\end{lemma}
\begin{proof}
Compare compact-support cohomology and homology Mayer--Vietoris by cap
product. Lift the boundary of a homology class through $D_{U\cap V}^2$,
then correct by classes from $U$ and $V$ to obtain surjectivity of
$D_{U\cup V}^1$. The next segment and the five lemma give the
degree-two isomorphism. Comparing cokernels gives degree three;
exactness preserves higher-degree vanishing.
\end{proof}

\begin{lemma}[The property on chart domains]
\label{lem:foundations-chart-cap}
Every open subset of a coordinate chart satisfies the cap property.
\lean{Poincare.Topology.integralOpenCapProperty_of_chart}
\uses{lem:foundations-cap-union}
\end{lemma}
\begin{proof}
Compute on an open box by excision from a cube modulo its boundary.
Use the union step for finite unions of boxes. Cycles, bounding chains,
and compact supports lie in finite stages, so the property passes to
directed unions and then through a chart.
\end{proof}

\begin{proposition}[Vanishing second homology]
\label{prop:foundations-second-homology}
A closed simply connected smooth three-manifold has $H_2(M)=0$.
\lean{Poincare.Topology.integralThreeManifoldHomologyTwo_isZero}
\uses{lem:foundations-chart-cap,lem:foundations-cap-union,prop:foundations-triangulation,lem:foundations-coherent}
\end{proposition}
\begin{proof}
Chart induction gives a surjection $H_c^1(M)\to H_2(M)$, and
compactness identifies the domain with $H^1(M)$. Integrating a singular
one-cocycle along paths from a basepoint gives a primitive: simple
connectedness makes the value independent of the path. Thus $H^1(M)=0$.
The formal proof transports this calculation through the finite
triangulation to accommodate arbitrary universes.
\end{proof}

\section{From homology to based homotopy}
\label{sec:foundations-hurewicz}

\begin{corollary}[The initial based classes]
\label{cor:foundations-homotopy-groups}
For every $p\in M$, $\pi_2(M,p)=0$ and $\pi_3(M,p)\cong\mathbb Z$.
In particular $\pi_3(M,p)$ contains a nonzero class.
\lean{Poincare.Topology.nonempty_threeManifoldTopologyConclusion_of_integralHomology_two_isZero}
\uses{prop:foundations-top-homology,prop:foundations-second-homology}
\end{corollary}
\begin{proof}
Integral Hurewicz in degree two identifies $\pi_2$ with $H_2=0$;
in degree three it then identifies $\pi_3$ with $H_3\cong\mathbb Z$.
Changing basepoint transports the groups by the isomorphism of a chosen
path. The class corresponding to $1$ remains nonzero.
\end{proof}

\section{The stronger sphere homotopy equivalence}
\label{sec:foundations-cw-equivalence}

\begin{proposition}[A three-dimensional cellular recognition result]
\label{prop:foundations-cw-equivalence}
Let $C$ be a nonempty path-connected Hausdorff CW complex of dimension
at most three. If $\pi_1(C,p)=\pi_2(C,p)=0$ and
$\pi_3(C,p)\cong\mathbb Z$ for every $p\in C$, then $C\simeq S^3$.
\lean{Poincare.Topology.nonempty_homotopyEquiv_sphere_of_cw_three_pi_int}
\end{proposition}
\begin{proof}
Choose a zero-cell $p$ and a map $f:S^3\to C$ representing a generator
of $\pi_3(C,p)$. Its map on $\pi_3$ is an isomorphism. Vanishing of
$\pi_1$ and $\pi_2$ contracts the inclusion of the two-skeleton
relative to $p$; extend this homotopy over $C$.

Each three-cell now gives a based sphere map. Lift its class through
$f_*$ and paste to obtain $g:C\to S^3$ with
$fg\simeq\mathrm{id}_C$. Injectivity of $f_*$ cancels $f$ from
$fgf\simeq f$, yielding $gf\simeq\mathrm{id}_{S^3}$.
\end{proof}

Apply this to the triangulated model and initial based classes. The
result is a homotopy equivalence; surgery reconstruction is needed for
a diffeomorphism.

\section{Excluding a two-sided projective plane}
\label{sec:foundations-projective-plane}

\begin{definition}[The excluded neighborhood]
\label{def:foundations-projective-neighborhood}
A space $X$ excludes a trivial normal projective-plane neighborhood if
there is no topological open embedding $\mathbb{RP}^2\times(-1,1)\to X$.
\lean{PoincareMT.NoTrivialNormalProjectivePlane}
\end{definition}

\begin{proposition}[Orientation exclusion]
\label{prop:foundations-projective-exclusion}
A Hausdorff second-countable smooth three-manifold with compatible
chart signs excludes a trivial normal projective-plane neighborhood.
Neither compactness nor simple connectedness is required.
\lean{PoincareMT.m83OrientationExclusion}
\uses{def:foundations-chart-signs,def:foundations-projective-neighborhood}
\source{MT2007}\dcref{Theorem 0.3, footnote 2, p. xii}
\end{proposition}
\begin{proof}
Correct negative chart signs by reflection to obtain a positive atlas.
A proposed open embedding pulls it back to
$\mathbb{RP}^2\times(-1,1)$ and supplies a local integral orientation;
no derivative of the embedding is used.

Pull this orientation to $S^2\times(-1,1)$. It is invariant under
$(u,t)\mapsto(-u,t)$. The shell identification $(u,t)\mapsto(2+t)u$
identifies the deck transformation with negation in $\mathbb R^3$,
which reverses local orientation. On the connected shell all orientations
differ by a constant sign, so invariance equates a generator with its
negative, a contradiction.
\end{proof}

\section{The forward geometric and backward topological arguments}
\label{sec:foundations-route}

Surgery uses the orientation and projective-plane exclusion. Filling
width uses a nonzero $\pi_3$ class on the fixed time-zero component,
transported to its chosen basepoint and converted to a class in the
actual $C^1$ free-loop space.
Section~\ref{sec:filling-fixed-initial-class} specifies those identifications.

The estimates force extinction of that same flow. Reading its finite
history backward reconstructs the connected-sum topology and reduces
the simply connected case to a sphere. Compatible smoothing finally
gives one homeomorphism from a topological manifold to a smooth model;
composition with the smooth endpoint proves the topological theorem.
